\documentclass{article}
\usepackage{amsmath,amssymb,amsthm}
\newcommand{\HH}{H_0^1(\Omega)}
\newtheorem{proposition}{Proposition}
\begin{document}
\section{Introduction}
\label{sec:introduction}
We studies the passage from weak to strong convergence.
This elementary example isolates the role of compactness.

\section{A claimed compactness argument}
Let $N\geq 3$, let $\Omega=B_1(0)\subset\mathbb R^N$, and set
$2^*=2N/(N-2)$, and fix $1\leq q<2^*$.
\begin{proposition}\label{prop:critical}
If $u_n\rightharpoonup u$ in $\HH$, then
$u_n\to u$ strongly in $L^q(\Omega)$.
\end{proposition}
\begin{proof}
The Sobolev embedding $\HH\hookrightarrow L^q(\Omega)$ is compact.
This prove the assertion.
\end{proof}

\section{A correct estimate}
Let $v_n,v\in L^2(\Omega)$ and assume $v_n\to v$ in $L^2(\Omega)$.
By the Cauchy--Schwarz inequality,
\[
 \|v_n^2-v^2\|_{L^1(\Omega)}
 \leq \|v_n-v\|_{L^2(\Omega)}
       \bigl(\|v_n\|_{L^2(\Omega)}+\|v\|_{L^2(\Omega)}\bigr)\to0.
\]
\end{document}
